\documentclass[11pt]{article}
\usepackage{setspace}
\usepackage[top=1in, bottom=1in, left=1in, right=1in]{geometry}
\usepackage{hyperref}
\onehalfspacing

% Title information
\title{Preferences on Tax Enforcement and Tax Administration:\\Note on Information Treatments}
\author{Tim Kircali, Stefanie Stantcheva, Matthias Weber}
\date{\today} % Use \date{} for no date

% Bibliography package
\usepackage[backend=biber, style=apa]{biblatex}
\addbibresource{references.bib} % Link to the bibliography file

\begin{document}
\maketitle
\section{General Strategy}
As we are using a commercial survey provider to conduct our survey, we need to be aware that respondents will complete these surveys wherever they are and may be distracted by their surroundings. Also, we can't and probably shouldn't force them to look at our information. People may not like being forced to look at it and may just do something else while we lock the screen. Instead, we should make our information appealing to look at, so that people might want to look at it. Then we can check how many have actually done so by asking very simple questions about the videos.

Below we suggest the information we would like to provide to respondents. Compared to other literature using survey experiments with information provision, the amount of information we provide is not small. Other papers often just update beliefs about a probability or a (relative place in a) distribution. Since we provide much more information, using textual information alone is not an option. Rather, we should use animated videos and definitely consider using audio as well, as it is much easier for respondents to process. The videos should be up to the video standards that people see in their daily lives, which have become very high in recent years. Videos should be mobile friendly, probably a TikTok/Reels/Youtube shorts format (1080x1920 pixels) would be appropriate. We could also consider implementing a slider to increase/decrease the video speed. This would allow us to use a default speed that is fairly low, but allow respondents who are used to higher video speeds to increase it. \\ \\
Videos that might be comparable: 
\begin{itemize}
    \item Videos in \cite{dechezlepretre2025fighting} include much information: They can be found here: \href{https://lse.eu.qualtrics.com/WRQualtricsControlPanel/File.php?F=F_3NNS6u7MbEm738y}{LINK}
    \item Videos in \cite{stantcheva2021understanding} do also include much information: They can be found here: \href{https://socialeconomicslab.org/research/publications/understanding-tax-policy-how-do-people-reason/}{LINK}
\end{itemize}
Videos that might be less comparable: 
\begin{itemize}
    \item Videos in \cite{alesina2023immigration} do not include much information, they are rather used for priming. They can be found here: \href{https://www.youtube.com/watch?v=_1SoLYX8OyE}{LINK}
    \item Another example that might be less suitable are informational treatments in \cite{alesina2018intergenerational}. The amount of information people receive is rather low, and video standards have improved since then (2016). They can be found here by clicking through the survey: \href{https://harvard.az1.qualtrics.com/jfe/form/SV_5dxninfErZ246X3}{LINK}
\end{itemize}

\section{Information Treatments}

\subsection{Tax Evasion Treatment}
\subsubsection*{General}
The first treatment provides respondents with information about the economic effects of third-party reporting and AEoI in tax enforcement. We give them information about the size of the tax gap and we directly point out that AEI is effective in reducing tax evasion by explaining evidence which shows that tax evasion is rare in areas which are covered under AEI. The primary goal of the tax evasion treatment is to manipulate people’s knowledge about the economic effects of tax enforcement and AEI without affecting considerations related to fairness \& inequality. \\ \\  
Estimated video length: 3.5-4 min (ChatGPT)
\subsubsection*{Information Text}
The following explanations could be used as subtitles and audio for an animated video.
\begin{itemize}
\item Tax revenue is essential for modern governments to provide their functions, reaching from the provision of safety over infrastructure, education and healthcare to social security.
\item One problem governments face is the loss of tax revenue due to tax evasion. This means that some taxpayers deliberately hide some of their income or wealth from the tax administration. 
\item The tax gap is the difference between the taxes that should be paid and the taxes that are actually paid. Several estimates show that this gap is very high in many countries.
\item In the US, for example, the tax administration, the Internal Revenue Service, estimates that only 85\% of taxes are paid voluntarily and on time. The remaining 15\% amount to \$696 billion and more than \$600 billion of these are never collected
\begin{itemize}
    \item Animation: These numbers will be visualized with diagrams.
\end{itemize}
\item Other countries have similar problems with tax evasion.
\begin{itemize}
    \item Animation: Show the tax gap in other countries (where available)
\end{itemize}
\item If the governments were able to collect these taxes, they  could use them for several purposes. To give a few examples, the US government could roughly double its spending on education (\$677 billion), it could increase social security by 50\% (\$1,219 billion), or the government could simply reduce taxes for all taxpayers by about 13\%. 
\begin{itemize}
    \item Animation: Show visually how strong government budget would change. 
\end{itemize}
\item Tax administrations have a number of options to combat tax evasion and therefore reduce the tax gap.
\item For example, they can increase their audit activities, which means that they investigate more tax returns or taxpayers in detail.  
\item Research has shown that tax evasion hardly takes place in areas that are covered by third-party reporting.
\item Third-party reporting means that certain institutions are obliged to transmit tax-relevant information to the tax administration directly.
\item For example, in many countries, including in XX, employers report to the tax administration how much they pay their employees in wages.
\item Another example is that if interest and dividend payments of assets held by a person in a bank are higher than a certain amount, the bank automatically informs the tax administration in some countries, including in XX.
\item Increasing the amount of data that is third-party reported is therefore an effective way to reduce tax evasion.
\begin{itemize}
    \item Animation: We can visualise two parts of the financial data, one that is reported to third parties and one that is not. We show that most of the tax gap is the result of the unreported part. We then show that the tax gap decreases as we increase the proportion that is third party reported.
\end{itemize}
\item And, recent research shows that not only domestic, but also international information exchange is very effective in reducing tax evasion. In 2017-2018, more than 100 countries started to cooperate internationally to combat tax evasion. Their tax administrations started to automatically exchange information with foreign banks to detect offshore tax evasion. In Denmark, the AEoI has led to a reduction in the offshore tax gap of about 70\%. 
\end{itemize}
\subsubsection*{Follow up questions}
We will ask 2-3 very basic questions to understand if respondents actually watched the video.
\subsection{Inequality Treatment}
\subsubsection*{General}
The second specific information treatment again includes information on tax evasion, but also explains the fact that the tax gap is much more pronounced for high-income individuals and discusses why this is the case. After that, we discuss fairness considerations of tax evasion, such as the fact that honest taxpayers have to finance the fiscal gap caused by tax evaders. The primary goal of the inequality treatment is to manipulate people’s knowledge about distributional consequences of third-party reporting and AEI.\\ \\  
Estimated video length: 3.5-4 min (ChatGPT)
\subsubsection*{Information Text}
The following explanations could be used as subtitles and audio for an animated video.
\begin{itemize}
    \item The tax gap is the difference between the taxes that should be paid and the taxes that are actually paid. Several estimates show that this gap is very high in many countries.
    \item In the US, for example, the tax administration, the Internal Revenue Service, estimates that only 85\% of taxes are paid voluntarily and on time. The remaining 15\% amount to \$696 billion and more than \$600 billion of these are never collected.
    \begin{itemize}
    \item Animation: These numbers will be visualized with diagrams.
    \end{itemize}
    \item Research has shown that not all groups of the population contribute equally to the tax gap. Taxpayers from low to medium income households only rarely evade taxes. Most of these taxpayers’ income stems from wages, pensions and investment income in domestic financial institutions. Such income is at least partially reported to the tax administration by a third party, so that taxpayers have few opportunities to evade taxes on these incomes. 
    \item Therefore, tax evasion mainly happens among wealthy taxpayers.
    \item A small group of the population is responsible for a large part of the tax gap. For example (x)\% of the absolute tax gap is due to the (y)\% richest part of the population. 
    \begin{itemize}
        \item Animation: We show how tax evasion is distributed along the income distribution and show how much each percentile contributes to the "tax evasion pie". 
        \item Here we could use examples of \cite{johns2010distribution}, which show evasion rates along the income distribution for deciles and quantiles.
    \end{itemize}  
    \item Recent research has shown that tax evasion is even unequally distributed among these mentioned (y)\%. An analysis of leaks of undeclared offshore bank accounts shows that almost all of them were held by people in the top 1\% of the income distribution. These findings suggest that only people at the very top of the income distribution are using offshore bank accounts to evade tax.
    \begin{itemize}
        \item Animation: We zoom in to the top decile of the income distribution and show that tax evasion is also highly unequally distributed at the top of the income distribution, with most of the evasion driven by the top 1\%.
        \item For that we can show evidence by \cite{alstadsaeter2019tax} where evasion rates are even higher than suggested in \cite{johns2010distribution}. 
    \end{itemize} 
\item The researchers also estimate that the super-rich, referring to the top 0.01 per cent of Scandinavian taxpayers, evade about 25 per cent of their tax liability just by hiding assets and investment income abroad. 
\begin{itemize}
    \item Animation: We show previously suggested evasion rates by quantile and show that the evasion estimate for top income by \cite{alstadsaeter2019tax} suggests even higher evasion rates. 
\end{itemize}
\item When a group of taxpayers evade taxes, the government has less money to spend on public services. It must either cut services or raise taxes, which will harm honest taxpayers who have paid their fair share of taxes. 
\begin{itemize}
    \item Animation: Show the government budget and the government budget gap. Animate how that affects honest taxpayers.
\end{itemize}
\item Tax administrations have a number of options to combat tax evasion and therefore reduce the tax gap.
\item For example, they can increase their audit activities, which means that they investigate more tax returns or taxpayers in detail.  
\item Research has shown that tax evasion hardly takes place in areas that are covered by third-party reporting.
\item Third-party reporting means that certain institutions are obliged to transmit tax-relevant information to the tax administration directly.
\item For example, in many countries, including in XX, employers report to the tax administration how much they pay their employees in wages.
\item Another example is that if interest and dividend payments of assets held by a person in a bank are higher than a certain amount, the bank automatically informs the tax administration in some countries, including in XX.
\item Increasing the amount of data that is third-party reported is therefore an effective way to reduce tax evasion.
\begin{itemize}
    \item Animation: We can visualise two parts of the financial data, one that is reported to third parties and one that is not. We show that most of the tax gap is the result of the unreported part. We then show that the tax gap decreases as we increase the proportion that is third party reported.
\end{itemize}
\item And, recent research shows that not only domestic, but also international information exchange is very effective in reducing tax evasion. In 2017-2018, more than 100 countries started to cooperate internationally to combat tax evasion. Their tax administrations started to automatically exchange information with foreign banks to detect offshore tax evasion. In Denmark, the AEoI has led to a reduction in the offshore tax gap of about 70\%. Recent research has shown that most of this reduction has been driven by the returns of the top 1\% of taxpayers.
\end{itemize}
\subsection{Data Privacy Treatment}
\subsubsection*{General}
In the data privacy treatment, we illustrate tax-relevant players and institutions and how tax-relevant data is being exchanged between them. We explain how information exchange would change if the level of third-party reporting were expanded. The primary goal of the data privacy treatment is to manipulate people’s knowledge about third-party reporting and AEoI, with focus on data privacy considerations.\\ \\  
Estimated video length: 3.5-4 min (ChatGPT)
\subsubsection*{Information Text}
The following explanations could be used as subtitles and audio for an animated video.
\begin{itemize}
    \item While 50 years ago collecting data about people was difficult, today – thanks to digital technology – it has become fast, cheap, and often invisible.
    \item Large companies can now track where people go using mobile phone data. This information is used not only to plan logistics, but also to profile individuals in ways they never really agreed to.
    \item They monitor what people search online and build detailed personal profiles to target them with ads – not just to match interests, but to influence decisions and behavior.
    \item The more data they collect, the more power they have to predict, steer, and even manipulate what people see, buy, or believe. Personalized content is not just about convenience – it’s about capturing attention and maximizing profit.
    \item Not only private companies but also governments increasingly collect data – sometimes to improve services, but also for surveillance and control.
    \item For example a government's tax administration does not only rely on tax payers self-reported information to evaluate tax returns. It also collects information reported by third-parties to check that the self-reported information is correct. 
    \item Third-party reporting means that certain institutions are obliged to transmit tax-relevant information to the tax administration directly.
    \item For example, in many countries, including in [Country XX], employers report to the tax administration how much they pay their employees in wages.
    \begin{itemize}
        \item We change the country we refer to depending on where we conduct the survey. 
    \end{itemize}
    \item Another example is that if interest and dividend payments of assets held by a person in a bank are higher than a certain amount, the bank automatically informs the tax administration in some countries, including in [Country XX]. 
    \begin{itemize}
        \item We change the country we refer to depending on where we conduct the survey.
    \end{itemize}
    \item (Here we might provide more detailed information on which data is being automatically exchanged)
    \item The level of third-party reporting, that is the amount of different data sources that is automatically reported to the tax administration, differs among countries. While in some countries nearly everything, from employer data to bank account and health insurance data is reported to the tax administration, in other countries tax administrations have hardly any automatic access to tax-relevant data.
    \begin{itemize}
        \item Animation: We highlight differences across countries by animating them. 
    \end{itemize}
    \item If the tax administration does not receive information automatically, it can either rely on the information provided by the taxpayer or it can take further action to obtain the information for a given tax payer. 
    \item Often, there are even governmental institutions that have tax-relevant data but do not report to the tax administration. In such situations, the government actually has the data but the tax administration has no access to it.
    \item Third-party reporting clearly helps tax administrations to directly verify the accuracy of tax reports. But some people are opposed to it because they are concerned about data privacy. They do not like the fact that governmental institutions collect and exchange data about them.
    \item On the one hand, it is correct that the tax administration has less data without third-party reporting.
    \item On the other, the tax administration has almost all tax-relevant data of all tax payers who file their tax reports honestly and correctly also without third-party reporting, because these data are then reported by the (honest) tax return filers directly. The tax administration only misses data of dishonest taxpayers and of those who do not need or want to file a tax return. Furthermore, the fact that tax administrations do not have the data does not mean that the data would not be available to governments if they turn non-benevolent if the data are already available at other governmental institutions (or to some extent at private companies, which can be forced by governments to hand over that data).
\end{itemize}
\subsubsection*{Follow up questions}
We will ask 2-3 very basic questions to understand if respondents actually watched the video.
\subsection{Tax Filing Treatment}
\subsubsection*{General}
The last treatment discusses concerns about tax filing costs such as the fact that an average American taxpayer has a tax filing effort and cost of about 13 hours + 270 \$. It describes how (guided) tax software works and what a prepopulated tax return looks like. We emphasize that tax software and prepopulated tax returns help tax filers to make fewer mistakes and significantly reduce their tax filing costs. Importantly, we highlight the importance of third-party reporting to the tax administration in providing prepopulated tax returns. The primary goal of the tax filing treatment is to manipulate people’s knowledge and perceptions about tax filing costs and tax filing services and how they are related to the existence of third-party reporting and AEI. \\ \\  
Estimated video length: 2.5-3 min (ChatGPT)
\subsubsection*{Information Text}
The following explanations could be used as subtitles and audio for an animated video.
\begin{itemize}
    \item Most people need to file tax returns to transmit information to the tax administration, such as received salary, bank account balances or paid contributions to health insurance.
    \item While some taxpayers prepare their tax returns themselves, others pay for advice or pay someone else to prepare their tax returns. 
    \item Taxpayers who prepare their tax returns themselves often use specific tax preparation software to file their taxes. Such software helps taxpayers file their tax returns by providing guidance through the tax formulas and explaining which deductions might apply. It provides field-by-field advice to reduce errors and foregone reductions.
    \begin{itemize}
        \item Animation: We will visualize how tax software works, such that people get an understanding of it. 
    \end{itemize}
    \item All these options are time consuming or costly. An estimate by the US tax administration, the Internal Revenue Service, suggests that taxpayers in the US need on average  13 hours plus \$270 for tax software and help to file their taxes.
    \item In some countries, the tax administrations offer solutions that reduce tax filing costs and efforts. The extent to which they do so differs among countries. While most taxpayers in the US have to pay for tax software licenses, many other countries offer free government-provided tax software, which makes it easier for taxpayers to file their taxes. 
    \begin{itemize}
        \item Animation: We make it visually clear that tax payers in some countries are provided with better resources, while others are not. 
    \end{itemize}
    \item Some countries even offer prepopulated tax returns. This means that the tax returns are pre-filled with information the tax administration already has. Prepopulated tax returns make the process of filing taxes much shorter and easier, with lower potential for making mistakes.
    \begin{itemize}
        \item Animation: We show what a prepopulated tax return looks like and clearly visualise that many/most fields are already prepopulated, thus reducing the taxpayer's filing burden. 
    \end{itemize}
    \item If the tax administration does not have the respective data, it will not be able to prepopulate much, and most fields will be left blank and have to be filled in manually by taxpayers. The more tax-relevant information the tax administration has, the more fields in the tax return can be prepopulated.
    \begin{itemize}
        \item Animation: We clearly illustrate that many fields remain empty when the level of information is low, and that taxpayers then have to file manually.
    \end{itemize}
    \item Tax administrations often get access to tax-relevant information by third-party reporting. Third-party reporting means that certain institutions are obliged to transmit tax-relevant information to the tax administration directly.
    \item The more information the tax administration automatically receives, the better it can provide pre-populating services. Therefore, it is not surprising that countries where many institutions automatically exchange information offer good prepopulating services for their taxpayers.
    \item For example, in Switzerland, where the level of third-party reporting is low, tax authorities can only prepopulate fields with outdated values from taxpayers’ previous tax returns, which taxpayers have to update. In other countries such as  Denmark and Netherlands most of the information is third-party reported, and therefore tax administrations in these countries are able to offer fully prepopulated tax returns to most taxpayers.
    \begin{itemize}
        \item Animation: Using colors, we can clearly visualise which fields are pre-populated with the current values, which are pre-populated with outdated values, and which are left blank.
    \end{itemize}
    \item (Potentially, we could discuss the following points.)
    \begin{itemize} 
        \item Even if government provided tax filing services such as free tax software and prepopulated tax returns bring many advantages, there are also people that argue against governments efforts to provide them.
        \item For example, they argue that tax software is already being offered by commercial providers and that government provided software might not reach its quality. 
        \item Another argument is that the IRS should not act as both the tax collector and the tax preparer. This could lead to conflicts of interest or undermine the interest representation of taxpayers.
    \end{itemize}
\end{itemize}
\subsubsection*{Follow up questions}
We will ask 2-3 very basic questions to understand if respondents actually watched the video.
\subsection{Full Info Treatment}
\subsubsection*{General}
The full information treatment includes information from all treatments. The tax evasion and inequality treatments contain a lot of overlapping information, so we present this information together in the first part to reduce repetition. After that, we present the contents of the data privacy and tax filing treatments, where we again reduce overlap.   
One problem is that the full information treatment is very long and people's attention might (systematically) vary during the video. For example, people might stop listening in the second part of the full info treatment. Or they might not remember the first part. To check for such effects, we ask simple validation questions after the full info treatment has been shown, each relating to the content of each specific treatment. We might consider doing a robustness check by changing the order in which the information is presented in the full info treatment.\\ \\  
Estimated video length: 9-13 min (ChatGPT)
\subsubsection*{Information Text}
\begin{itemize}
    \item Tax revenue is essential for modern governments to provide their functions, reaching from the provision of safety over infrastructure, education and healthcare to social security.
    \item One problem governments face is the loss of tax revenue due to tax evasion. This means that some taxpayers deliberately hide some of their income or wealth from the tax administration. 
    \item The tax gap is the difference between the taxes that should be paid and the taxes that are actually paid. Several estimates show that this gap is very high in many countries.
    \item In the US, for example, the tax administration, the Internal Revenue Service, estimates that only 85\% of taxes are paid voluntarily and on time. The remaining 15\% amount to \$696 billion and more than \$600 billion of these are never collected.
    \begin{itemize}
    \item Animation: These numbers will be visualized with diagrams.
    \end{itemize}
    \item Other countries have similar problems with tax evasion.
    \begin{itemize}
    \item Animation: Show the tax gap in other countries (where available)
    \end{itemize}
    \item Research has shown that not all groups of the population contribute equally to the tax gap. Taxpayers from low to medium income households only rarely evade taxes. Most of these taxpayers’ income stems from wages, pensions and investment income in domestic financial institutions. Such income is at least partially reported to the tax administration by a third party, so that taxpayers have few opportunities to evade taxes on these incomes. 
    \item Therefore, tax evasion mainly happens among wealthy taxpayers.
    \item A small group of the population is responsible for a large part of the tax gap. For example (x)\% of the absolute tax gap is due to the (y)\% richest part of the population. 
    \begin{itemize}
        \item Animation: We show how tax evasion is distributed along the income distribution and show how much each percentile contributes to the "tax evasion pie". 
        \item Here we could use examples of \cite{johns2010distribution}, which show evasion rates along the income distribution for deciles and quantiles.
    \end{itemize}  
    \item Recent research has shown that tax evasion is even unequally distributed among these mentioned (y)\%. An analysis of leaks of undeclared offshore bank accounts shows that almost all of them were held by people in the top 1\% of the income distribution. These findings suggest that only people at the very top of the income distribution are using offshore bank accounts to evade tax.
    \begin{itemize}
        \item Animation: We zoom in to the top decile of the income distribution and show that tax evasion is also highly unequally distributed at the top of the income distribution, with most of the evasion driven by the top 1\%.
        \item For that we can show evidence by \cite{alstadsaeter2019tax} where evasion rates are even higher than suggested in \cite{johns2010distribution}. 
    \end{itemize} 
\item The researchers also estimate that the super-rich, referring to the top 0.01 per cent of Scandinavian taxpayers, evade about 25 per cent of their tax liability just by hiding assets and investment income abroad. 
\begin{itemize}
    \item Animation: We show previously suggested evasion rates by quantile and show that the evasion estimate for top income by \cite{alstadsaeter2019tax} suggests even higher evasion rates. 
\end{itemize}
\item When a group of taxpayers evade taxes, the government has less money to spend on public services. It must either cut services or raise taxes, which will harm honest taxpayers who have paid their fair share of taxes. 
\item In fact if the governments were able to collect the evaded taxes, they could use them for several purposes. To give a few examples, the US government could roughly double its spending on education (\$677 billion), it could increase social security by 50\% (\$1,219 billion), or the government could simply reduce taxes for all taxpayers by about 13\%.
\begin{itemize}
    \item Animation: Show visually how strong government budget would change. 
\end{itemize}
\item Tax administrations have a number of options to combat tax evasion and therefore reduce the tax gap.
\item For example, they can increase their audit activities, which means that they investigate more tax returns or taxpayers in detail.  
\item Research has shown that tax evasion hardly takes place in areas that are covered by third-party reporting.
\item Third-party reporting means that certain institutions are obliged to transmit tax-relevant information to the tax administration directly.
\item For example, in many countries, including in XX, employers report to the tax administration how much they pay their employees in wages.
\item Another example is that if interest and dividend payments of assets held by a person in a bank are higher than a certain amount, the bank automatically informs the tax administration in some countries, including in XX.
\item Increasing the amount of data that is third-party reported is therefore an effective way to reduce tax evasion.
\begin{itemize}
    \item Animation: We can visualise two parts of the financial data, one that is reported to third parties and one that is not. We show that most of the tax gap is the result of the unreported part. We then show that the tax gap decreases as we increase the proportion that is third party reported.
\end{itemize}
\item And, recent research shows that not only domestic, but also international information exchange is very effective in reducing tax evasion. In 2017-2018, more than 100 countries started to cooperate internationally to combat tax evasion. Their tax administrations started to automatically exchange information with foreign banks to detect offshore tax evasion. In Denmark, the AEoI has led to a reduction in the offshore tax gap of about 70\%. Recent research has shown that most of this reduction has been driven by the returns of the top 1\% of taxpayers.
    \item The advances in data collection, that make domestic and international exchange of information possible is mainly result of digital technology. While 50 years ago collecting data about people was difficult, today it has become fast, cheap, and often invisible.
    \item Large companies can now track where people go using mobile phone data. This information is used not only to plan logistics, but also to profile individuals in ways they never really agreed to.
    \item They monitor what people search online and build detailed personal profiles to target them with ads – not just to match interests, but to influence decisions and behavior.
    \item The more data they collect, the more power they have to predict, steer, and even manipulate what people see, buy, or believe. Personalized content is not just about convenience – it’s about capturing attention and maximizing profit.
    \item Not only private companies but also governments increasingly collect data – sometimes to improve services, but also for surveillance and control.
    \item As we have mentioned before, a government's tax administration also collects information reported by third-parties to check that the self-reported information is correct. 
    \item (Here we might provide more detailed information on which data is being automatically exchanged)
    \item The level of third-party reporting, that is the amount of different data sources that is automatically reported to the tax administration, differs among countries. While in some countries nearly everything, from employer data to bank account and health insurance data is reported to the tax administration, in other countries tax administrations have hardly any automatic access to tax-relevant data.
    \begin{itemize}
        \item Animation: We highlight differences across countries by animating them. 
    \end{itemize}
    \item If the tax administration does not receive information automatically, it can either rely on the information provided by the taxpayer or it can take further action to obtain the information for a given tax payer. 
    \item Often, there are even governmental institutions that have tax-relevant data but do not report to the tax administration. In such situations, the government actually has the data but the tax administration has no access to it.
    \item Third-party reporting clearly helps tax administrations to directly verify the accuracy of tax reports. But some people are opposed to it because they are concerned about data privacy. They do not like the fact that governmental institutions collect and exchange data about them.
    \item On the one hand, it is correct that the tax administration has less data without third-party reporting.
    \item On the other, the tax administration has almost all tax-relevant data of all tax payers who file their tax reports honestly and correctly also without third-party reporting, because these data are then reported by the (honest) tax return filers directly. The tax administration only misses data of dishonest taxpayers and of those who do not need or want to file a tax return. Furthermore, the fact that tax administrations do not have the data does not mean that the data would not be available to governments if they turn non-benevolent if the data are already available at other governmental institutions (or to some extent at private companies, which can be forced by governments to hand over that data).
    \item Another interesting argument in favor of comprehensive third party reporting to the tax administration is that it can help reduce tax filling costs. 
    \item Most people need to file tax returns to transmit information to the tax administration, such as received salary, bank account balances or paid contributions to health insurance.
    \item While some taxpayers prepare their tax returns themselves, others pay for tax software licences, advice or pay someone else to prepare their tax returns. 
    \item In any case, all these options are time consuming or costly. An estimate by the US tax administration, the Internal Revenue Service, suggests that taxpayers in the US need on average  13 hours plus \$270 for tax software and help to file their taxes.
    \item In some countries, the tax administrations offer solutions that reduce tax filing costs and efforts. The extent to which they do so differs among countries. While most taxpayers in the US have to pay for tax software licenses, many other countries offer free government-provided tax software. Such software helps taxpayers file their tax returns by providing guidance through the tax formulas and explaining which deductions might apply. It provides field-by-field advice to reduce errors and foregone reductions.
    \begin{itemize}
        \item Animation: We make it visually clear that tax payers in some countries are provided with better resources, while others are not. 
    \end{itemize}
    \item Some countries even offer prepopulated tax returns. This means that the tax returns are pre-filled with information the tax administration already has. Prepopulated tax returns make the process of filing taxes much shorter and easier, with lower potential for making mistakes.
    \begin{itemize}
        \item Animation: We show what a prepopulated tax return looks like and clearly visualise that many/most fields are already prepopulated, thus reducing the taxpayer's filing burden. 
    \end{itemize}
    \item If the tax administration does not have the respective data, it will not be able to prepopulate much, and most fields will be left blank and have to be filled in manually by taxpayers. The more tax-relevant information the tax administration has, the more fields in the tax return can be prepopulated.
    \begin{itemize}
        \item Animation: We clearly illustrate that many fields remain empty when the level of information is low, and that taxpayers then have to file manually.
    \end{itemize}
    %\item Tax administrations often get access to tax-relevant information by third-party reporting. Third-party reporting means that certain institutions are obliged to transmit tax-relevant information to the tax administration directly.
    \item The more information the tax administration automatically receives, the better it can provide pre-populating services. Therefore, it is not surprising that countries where many institutions automatically exchange information offer good prepopulating services for their taxpayers.
    \item For example, in Switzerland, where the level of third-party reporting is low, tax authorities can only prepopulate fields with outdated values from taxpayers’ previous tax returns, which taxpayers have to update. In other countries such as  Denmark and Netherlands most of the information is third-party reported, and therefore tax administrations in these countries are able to offer fully prepopulated tax returns to most taxpayers.
    \begin{itemize}
        \item Animation: Using colors, we can clearly visualise which fields are pre-populated with the current values, which are pre-populated with outdated values, and which are left blank.
    \end{itemize}
    \item (Potentially, we could discuss the following points.)
    \begin{itemize} 
        \item Even if government provided tax filing services such as free tax software and prepopulated tax returns bring many advantages, there are also people that argue against governments efforts to provide them.
        \item For example, they argue that tax software is already being offered by commercial providers and that government provided software might not reach its quality. 
        \item Another argument is that the IRS should not act as both the tax collector and the tax preparer. This could lead to conflicts of interest or undermine the interest representation of taxpayers.
    \end{itemize}
\end{itemize}
\subsubsection*{Follow up questions}
We will ask 4-8 very basic questions to understand if respondents actually watched all parts of the video.

\newpage
\printbibliography

\end{document}